
\documentclass{article}
\usepackage{colortbl}
\usepackage{makecell}
\usepackage{multirow}
\usepackage{supertabular}

\begin{document}

\newcounter{utterance}

\centering \large Interaction Transcript for game `cladder', experiment `full\_v1.5\_default', episode 3946 with selfrefine{-}Qwen3.5{-}27B.
\vspace{24pt}

{ \footnotesize  \setcounter{utterance}{1}
\setlength{\tabcolsep}{0pt}
\begin{supertabular}{c@{$\;$}|p{.15\linewidth}@{}p{.15\linewidth}p{.15\linewidth}p{.15\linewidth}p{.15\linewidth}p{.15\linewidth}}
    \# & \multicolumn{2}{c}{Player} && \multicolumn{2}{c}{Game Master} \\
    \hline

    \theutterance \stepcounter{utterance}  
    & & & \multicolumn{4}{p{0.6\linewidth}}{
        \cellcolor[rgb]{0.9,0.9,0.9}{
            \makecell[{{p{\linewidth}}}]{
                \texttt{\tiny{[P1$\langle$GM]}}
                \texttt{You are an expert in causal inference. The following question is not a typical commonsense query, but rather a meticulously designed question created by a professor specializing in causal inference, intended to assess the students' mastery of the course content.} \\
\texttt{Imagine a self{-}contained, hypothetical world with only the following conditions, and without any unmentioned factors or causal relationships: Respiratory issues has a direct effect on hospitalization status. Black hair has a direct effect on hospitalization status. For hospitalized individuals, the correlation between respiratory issues and black hair is {-}0.08. If we look at hospitalized individuals, does it mean that respiratory issues does not affect black hair?} \\
\texttt{Start your answer with 'Yes' or 'No', followed by additional reasoning or evidence to support your explanation.} \\
            }
        }
    }
     \\ \\

    \theutterance \stepcounter{utterance}  
    & & \multicolumn{4}{p{0.6\linewidth}}{
        \cellcolor[rgb]{0.95,0.95,0.95}{
            \makecell[{{p{\linewidth}}}]{
                \texttt{\tiny{[P1$\leftrightarrow$P]}}
                \texttt{{[}DRAFT{]}} \\
\texttt{No} \\
\\ 
\texttt{In this hypothetical world, the causal structure is defined such that "Respiratory issues" and "Black hair" are both direct causes of "Hospitalization status," but there is no causal link between "Respiratory issues" and "Black hair" themselves. This setup creates a classic **collider** structure: \$Respiratory \textbackslash{}rightarrow Hospitalization \textbackslash{}leftarrow BlackHair\$.} \\
\\ 
\texttt{According to the principles of causal inference, conditioning on a collider (in this case, restricting the analysis to only "hospitalized individuals") induces a non{-}causal statistical association (spurious correlation) between its causes, even if they are independent in the general population. This phenomenon is known as **collider bias** or **Berkson's paradox**.} \\
\\ 
\texttt{The observed correlation of {-}0.08 among hospitalized individuals is a mathematical artifact of the selection process, not evidence of a causal relationship. The fact that a correlation exists (or does not exist) in the selected subgroup does not determine the causal direction or existence between the two variables. Since the problem statement explicitly defines the world as having *only* the mentioned causal relationships (where neither variable affects the other), the correlation observed in the hospitalized group is purely a result of conditioning on the common effect. Therefore, the presence of this correlation does not imply that respiratory issues affect black hair; rather, it confirms the operation of collider bias in a system where the two variables are causally independent.} \\
            }
        }
    }
    & \\ \\

    \theutterance \stepcounter{utterance}  
    & & \multicolumn{4}{p{0.6\linewidth}}{
        \cellcolor[rgb]{0.95,0.95,0.95}{
            \makecell[{{p{\linewidth}}}]{
                \texttt{\tiny{[P1$\leftrightarrow$P]}}
                \texttt{{[}SELF{-}FEEDBACK{]}} \\
\texttt{FORMAT: ok} \\
\texttt{RULES: ok} \\
\texttt{STRATEGY: ok} \\
\texttt{REASON: The answer correctly starts with "No" as required, accurately applies causal inference principles regarding collider bias to the hypothetical scenario, and provides sound reasoning without violating any constraints.} \\
            }
        }
    }
    & \\ \\

    \theutterance \stepcounter{utterance}  
    & \multicolumn{4}{p{0.6\linewidth}}{
        \cellcolor[rgb]{0.561,0.737,0.561}{
            \makecell[{{p{\linewidth}}}]{
                \texttt{\tiny{[P1$\rangle$GM]}}
                \texttt{No} \\
\\ 
\texttt{In this hypothetical world, the causal structure is defined such that "Respiratory issues" and "Black hair" are both direct causes of "Hospitalization status," but there is no causal link between "Respiratory issues" and "Black hair" themselves. This setup creates a classic **collider** structure: \$Respiratory \textbackslash{}rightarrow Hospitalization \textbackslash{}leftarrow BlackHair\$.} \\
\\ 
\texttt{According to the principles of causal inference, conditioning on a collider (in this case, restricting the analysis to only "hospitalized individuals") induces a non{-}causal statistical association (spurious correlation) between its causes, even if they are independent in the general population. This phenomenon is known as **collider bias** or **Berkson's paradox**.} \\
\\ 
\texttt{The observed correlation of {-}0.08 among hospitalized individuals is a mathematical artifact of the selection process, not evidence of a causal relationship. The fact that a correlation exists (or does not exist) in the selected subgroup does not determine the causal direction or existence between the two variables. Since the problem statement explicitly defines the world as having *only* the mentioned causal relationships (where neither variable affects the other), the correlation observed in the hospitalized group is purely a result of conditioning on the common effect. Therefore, the presence of this correlation does not imply that respiratory issues affect black hair; rather, it confirms the operation of collider bias in a system where the two variables are causally independent.} \\
            }
        }
    }
    & & \\ \\

    \theutterance \stepcounter{utterance}  
    & & & \multicolumn{2}{p{0.3\linewidth}}{
        \cellcolor[rgb]{0.9,0.9,0.9}{
            \makecell[{{p{\linewidth}}}]{
                \texttt{\tiny{[GM$|$GM]}}
                \texttt{No} \\
            }
        }
    }
    & & \\ \\

    \theutterance \stepcounter{utterance}  
    & & & \multicolumn{2}{p{0.3\linewidth}}{
        \cellcolor[rgb]{0.9,0.9,0.9}{
            \makecell[{{p{\linewidth}}}]{
                \texttt{\tiny{[GM$|$GM]}}
                \texttt{yes} \\
            }
        }
    }
    & & \\ \\

    \theutterance \stepcounter{utterance}  
    & & & \multicolumn{2}{p{0.3\linewidth}}{
        \cellcolor[rgb]{0.9,0.9,0.9}{
            \makecell[{{p{\linewidth}}}]{
                \texttt{\tiny{[GM$|$GM]}}
                \texttt{game\_result = LOSE} \\
            }
        }
    }
    & & \\ \\

\end{supertabular}
}

\end{document}
